% Part:sets-functions-relations
% Chapter: sets
% Section: schroder-bernstein

\documentclass[../../../include/open-logic-section]{subfiles}

\begin{document}

\olfileid[fr]{sfr}{siz}{sb}

\olsection{La notion de taille et le théorème de Schr\"oder-Bernstein}

\begin{explain}
Voici une idée intuitive : si $A$ n'est pas plus grand que $B$ et
si $B$ n'est pas plus grand que $A$, alors $A$ et $B$ sont
équipotents. À vrai dire, si cette idée était \emph{fausse}, nous
aurions bien du mal à justifier que notre définition de l'équipotence
exprime une comparaison des «~tailles~» des ensembles !{} Heureusement,
cette intuition est correcte, comme l'établit le théorème de
Schr\"oder-Bernstein.
\end{explain}

\begin{thm}[Schr\"oder-Bernstein]
  \ollabel{thm:schroder-bernstein}
  Si $\cardle{A}{B}$ et $\cardle{B}{A}$, alors $\cardeq{A}{B}$.
\end{thm}

\begin{explain}
Autrement dit, s'il existe une !!{injection} de $A$ dans~$B$ et
une !!{injection} de $B$ dans~$A$, alors il existe une
!!{bijection} de $A$ dans~$B$.

Ce résultat est toutefois assez \emph{difficile} à démontrer. Cantor
l'avait énoncé, mais ce sont d'autres mathématiciens qui l'ont
démontré.\footnote{Pour plus de détails historiques, voir par exemple
\citet[pp.~165--6]{Potter2004}.}
\oliflabeldef{sfr:cardinals:card-sb:sec}{Nous ne disposerons des
outils nécessaires pour \emph{démontrer} le théorème de
Schr\"oder-Bernstein qu'à la
\olref[sfr][cardinals][card-sb]{sec}.}{}%
Pour le moment, vous pouvez, et devez, l'admettre.

Le théorème de Schr\"oder-Bernstein est heureusement \emph{vrai} :
il justifie que nous considérions les relations $\cardeq{A}{B}$
et $\cardle{A}{B}$ comme des comparaisons de « taille ». Il est
aussi très \emph{utile}. Construire une !!{bijection} entre deux
ensembles équipotents peut être difficile. Le théorème de
Schr\"oder-Bernstein permet de scinder ce travail en deux :
il suffit de construire une !!{injection} du premier ensemble
dans le second, puis une autre dans le sens inverse.
\end{explain}
\end{document}
